% !PRINTED 220
\weberpage{220}
%<<CONT
octahedral substitutions absolutely unchanged, while $K$ again
suffers the same changes of sign as $f$.

\begin{center}
\S.\,58.

\gesperrt{The icosahedral group.}
\end{center}

Since in the icosahedral group we have only one system of conjugate
five-fold poles, we can (\S.\,53, 1.) assume in this group
the two substitutions

\weq{1}{\Theta(x) = \varepsilon x, \quad \psi(x) = \frac{-1}{x}}

where $\varepsilon$ denotes a primitive fifth root of unity.
We here take, as we are free to do, the substitution $\psi$ in the
form $-1:x$, because the formulae thereby become simpler.
The fundamental form of degree 12 belonging to the system of the
five-fold poles must, since it admits the substitutions (1), be of
the form:

\weq{2}{f(x_1,x_2) = x_1x_2(x_1^{10} + m\,x_1^5x_2^5 - x_2^{10}),}

and it remains to determine the constant
factor $m$. The two other fundamental forms are of degree 20
and degree 30. We can now determine $m$ by the theorem \S.\,54, 4.
if we can form a covariant of $f(x)$
whose degree is lower than 60, and which cannot be represented
as a product of powers of three functions of degree 12, 20, 30,
and which by the theorem mentioned must be identically
zero.

Now a covariant of degree 16 of the
form $f$ can easily be formed, if, following Bd.\ I, \S.\,60, we take the fourth polar of the
form $f(x_1,x_2)$:

\[
12\,.\,11\,.\,10\,.\,9\;P_4(x,\xi) = u_0\xi_1^4 + u_1\xi_1^3\xi_2 + u_2\xi_1^2\xi_2^2 + u_3\xi_1\xi_2^3 + u_4\xi_2^4,
\]

where

\[
u_0 = \frac{\partial^4 f}{\partial x_1^4}, \quad u_1 = 4\,\frac{\partial^4 f}{\partial x_1^3\,\partial x_2}, \quad u_2 = 6\,\frac{\partial^4 f}{\partial x_1^2\,\partial x_2^2},
\]
\[
u_3 = 4\,\frac{\partial^4 f}{\partial x_1\,\partial x_2^3}, \quad u_4 = \frac{\partial^4 f}{\partial x_2^4}.
\]

From this we obtain a covariant of degree 16 of $f$ as
%CONT>>
